%
% Paper title
% Authors
% 
%-----------------------------------------------------------------

\documentclass[12pt,longbibliography]{article}

% Packages
%-----------------------------------------------------------------

% Allow direct use of accents such as á é ñ.
\usepackage[utf8]{inputenc}

% This is a good idea to have some symbols included within the font
% properly:
% http://tex.stackexchange.com/questions/664/why-should-i-use-usepackaget1fontenc
\usepackage[T1]{fontenc}

% Set type of paper and margins.
\usepackage[a4paper, margin=2.7cm]{geometry}

\usepackage{amsmath, amsthm, amsfonts, amssymb}
\usepackage{mathrsfs}           % \mathscr font.

% Generate a PDF with hyperlinks in references.
\usepackage[colorlinks=true,linkcolor=black,citecolor=blue,urlcolor=black,breaklinks]{hyperref}

% Double-stroke font (\mathbbm).
\usepackage{bbm} 

\usepackage[spanish]{babel}
\usepackage{graphicx}
\usepackage{listings}
\usepackage{listings}
\usepackage{xcolor}

\definecolor{dartKeyword}{RGB}{0,119,170}
\definecolor{dartString}{RGB}{105,155,53}
\definecolor{dartComment}{RGB}{112,128,144}
\definecolor{dartType}{RGB}{38,139,210}

\lstset{
	language=Java,
	basicstyle=\ttfamily\small,
	keywordstyle=\color{dartKeyword}\bfseries,
	stringstyle=\color{dartString},
	commentstyle=\color{dartComment}\itshape,
	morekeywords={class, static, final, async, await, extends, abstract, override, Widget, dynamic, var, return},
	frame=single,
	rulecolor=\color{gray!30},
	backgroundcolor=\color{gray!5},
	breaklines=true,
	showstringspaces=false,
	tabsize=2,
	captionpos=b
}

% Bibliography
%-----------------------------------------------------------------

% This uses a bibliography style which hyperlinks the paper titles to
% the paper URL specified in the bibtex file. It also uses natbib,
% which cites papers by name such as Euler (1770) instead of [17].

\usepackage{breakurl}
\usepackage{natbib}
\usepackage{url}
\bibliographystyle{plainnat-linked}
% \bibliographystyle{plain}

% Shortcuts
%-----------------------------------------------------------------

% Absolute value 
\newcommand{\abs}[1]{\left\vert#1\right\vert}
% Inner product
\newcommand{\ap}[1]{\left\langle#1\right\rangle}
% Norm
\newcommand{\norm}[1]{\left\Vert#1\right\Vert}
% Triple norm
\newcommand{\tnorm}[1]{\left\vert\!\left\vert\!\left\vert#1\right\vert\!\right\vert\!\right\vert}

% Common double-stroke letters
\def\C{\mathbb{C}}
\def\N{\mathbb{N}}
\def\R{\mathbb{R}}
\def\S{\mathbb{S}}
\def\Z{\mathbb{Z}}

\def\D{\mathcal{D}}
% \def\L{\mathcal{L}}

\def\grad{\nabla}
\DeclareMathOperator{\dive}{div}
\DeclareMathOperator{\supp}{supp}
\DeclareMathOperator{\essup}{ess\,sup}
\DeclareMathOperator{\Lip}{Lip}
\DeclareMathOperator{\sgn}{sgn}

% Notation for differentials
\def\d{\,\mathrm{d}}
\def\dv{\d v}
\def \ddt{\frac{\mathrm{d}}{\mathrm{d}t}}
\def \ddt{\frac{\mathrm{d}}{\mathrm{d}t}}
\def \ddr{\frac{\mathrm{d}}{\mathrm{d}r}}

\def\p{\partial}

\def\ird{\int_{\R^d}}
\renewcommand{\L}[1]{L^{#1}(\R^d)}

% Theorems
%-----------------------------------------------------------------
\newtheorem{thm}{Theorem}[section]
% \newtheorem{thm}{Theorem}
\newtheorem{cor}[thm]{Corollary}
\newtheorem{lem}[thm]{Lemma}
\newtheorem{prp}[thm]{Proposition}
% \newtheorem{hyp}[thm]{Hypothesis}
\newtheorem{hyp}{Hypothesis}
\newtheorem{prb}{Problem}
\theoremstyle{definition}
\newtheorem{dfn}[thm]{Definition}
\theoremstyle{remark}
\newtheorem{rem}[thm]{Remark}
\theoremstyle{example}
\newtheorem{ex}[thm]{Example}


% Title, author, date
%-----------------------------------------------------------------

% If set, these will be the internal title and author of the PDF (and
% will be listed for example in ereaders and tablets).

% \hypersetup{pdftitle={Title of the PDF}}
% \hypersetup{pdfauthor={Author of the PDF}}


% Paper title and author
\title{REPORTE:\\Implementación de patrones estructurales en una aplicacion movil desarrollada en Flutter}
\author{Luisa Fernanda Hernández Hdez.}

% Date is set automatically unless specified.
% \date{October 2015}


% Addresses. The "article" class does not have a standard way to
% include them. This imitates the behaviour of "amsart" by inserting
% them at the end of the paper.
%-----------------------------------------------------------------

%


% =========================================================================

\begin{document}
	
	\maketitle
	\begin{center}
		\includegraphics[scale=0.3]{/Users/Admin/Downloads/portada.jpg}
	\end{center}
	\vfil
	\begin{center}{\large 
			Universidad Veracruzana\\Ingeniería de Software\\}
	\end{center}
%	\begin{abstract}
		
%	\end{abstract}
	\vspace{3cm}
	\pagebreak
	
	\tableofcontents
	
	\pagebreak
	
	\section{Introducción}
		Los patrones de diseño estructurales se enfocan en la composición de clases y objetos para formar estructuras de software más complejas y eficientes. Su objetivo primordial es facilitar que diferentes partes de un sistema, a menudo con interfaces incompatibles, puedan trabajar juntas de manera armónica sin comprometer la integridad del sistema.
		\\
		
		A diferencia de los patrones creacionales, que gestionan el proceso de instanciación, los estructurales utilizan la herencia y la composición para simplificar el diseño y reducir el acoplamiento. Permiten que los desarrolladores identifiquen una forma sencilla de realizar relaciones entre entidades, lo que resulta en sistemas más fáciles de mantener y escalar.
		
	\section{Descripción del Proyecto: Aplicación en Flutter y API REST}
		El sistema desarrollado representa una solución integral para la gestión de interacciones sociales (tweets, likes y comentarios). La arquitectura se divide en dos componentes principales:
		
		\begin{enumerate}
			\item \textbf{Backend (API REST):} Construido sobre el ecosistema de Java con Spring Boot. Utiliza una base de datos PostgreSQL y seguridad basada en JWT. El despliegue se realizó mediante contenedores Docker, garantizando un entorno de ejecución aislado y escalable en la nube (Render).
			\item \textbf{Frontend (Flutter):} Una aplicación móvil y web desarrollada bajo el paradigma de programación reactiva. Esta aplicación consume los recursos de la API de forma asíncrona, permitiendo la visualización de datos dinámicos y la persistencia de interacciones del usuario en tiempo real.
		\end{enumerate}
		
	\section{Implementación de Patrones Estructurales en la aplicación de Flutter}
		En Flutter, donde la composición es el principio rector (todo es un \textit{Widget}), los patrones estructurales permiten organizar la UI y la lógica de negocio de manera profesional.
	
		\subsection{Patrón Adapter}
		Se utiliza para transformar los datos crudos (JSON) de la API de Spring Boot en objetos de modelo de Dart.
		\begin{lstlisting}[language=dart, caption=Ejemplo de Adapter para Modelos]
			class TweetAdapter {
				static Tweet fromJson(Map<String, dynamic> json) {
					return Tweet(
					id: json['id'],
					content: json['text_content'], // Adapta 'text_content' a 'content'
					author: json['user']['name'],
					);
				}
			}
		\end{lstlisting}
		
		\subsection{Patrón Bridge}
		Separa la abstracción del servicio de su implementación específica por plataforma.
		\begin{lstlisting}[language=dart, caption=Ejemplo de Bridge para Notificaciones]
			abstract class NotificationImpl {
				void send(String msg);
			}
			
			class AndroidNotification extends NotificationImpl {
				@override
				void send(String msg) => print("Push Android: $msg");
			}
			
			class WebNotification extends NotificationImpl {
				@override
				void send(String msg) => print("Toast Web: $msg");
			}
		\end{lstlisting}
		
		\subsection{Patrón Composite}
		Es la base de Flutter; permite tratar widgets simples y complejos de forma uniforme.
		\begin{lstlisting}[language=dart, caption=Uso de Composite en UI]
			Widget build(BuildContext context) {
				return Column( // Widget compuesto
				children: [
				Text("Tweet de Luisa"), // Widget simple
				Icon(Icons.favorite),   // Widget simple
				],
				);
			}
		\end{lstlisting}
		
		\subsection{Patrón Decorator}
		Añade responsabilidades adicionales a un widget envolviéndolo en otro (Wrapper).
		\begin{lstlisting}[language=dart, caption=Ejemplo de Decorator con Container]
			Widget decoratedImage(Widget image) {
				return Container(
				decoration: BoxDecoration(
				border: Border.all(color: Colors.blue, width: 2),
				borderRadius: BorderRadius.circular(10),
				),
				child: image, // El widget original es "decorado"
				);
			}
		\end{lstlisting}
		
		\subsection{Patrón Facade}
		Proporciona una interfaz simplificada para el subsistema de red (API de Render).
		\begin{lstlisting}[language=dart, caption=Facade para Servicios API]
			class InteractionFacade {
				final ApiService _api = ApiService();
				
				Future<void> likeTweet(int id) async {
					await _api.post('/interactions/react/$id', body: 1);
				}
			}
		\end{lstlisting}
		
		\subsection{Patrón Flyweight}
		Optimiza la memoria compartiendo instancias de objetos pesados o repetitivos.
		\begin{lstlisting}[language=dart, caption=Uso de Flyweight en Estilos]
			class TextStyleFactory {
				static const TextStyle tweetStyle = TextStyle(
				fontSize: 14,
				fontWeight: FontWeight.bold,
				color: Colors.black,
				);
				// Se reutiliza la misma instancia en todos los tweets
			}
		\end{lstlisting}
		
		\subsection{Patrón Proxy}
		Actúa como un sustituto de un objeto para controlar el acceso o mostrar un estado de carga.
		\begin{lstlisting}[language=dart, caption=Proxy para Imágenes de Red]
			Image.network(
			'https://render-api.com/image.png',
			loadingBuilder: (context, child, progress) {
				if (progress == null) return child;
				return CircularProgressIndicator(); // Proxy visual mientras carga
			},
			);
		\end{lstlisting}
	
	\section{Conclusión: Importancia y Beneficios}
		La adopción de patrones estructurales trasciende la simple organización del código; es una garantía de calidad técnica. Sus beneficios principales incluyen:
		\begin{itemize}
			\item \textbf{Flexibilidad:} Permiten realizar cambios profundos en una 	parte del sistema sin afectar el resto.
			\item \textbf{Reutilización:} Facilitan el uso de componentes existentes en 	nuevas estructuras.
			\item \textbf{Legibilidad:} Establecen un estándar de comunicación claro para el equipo de desarrollo.
		\end{itemize}
	
		En conclusión, la correcta estructuración del software asegura que el proyecto sea resiliente, capaz de evolucionar ante nuevos requerimientos y técnicamente sostenible a largo plazo.
	
	
	\end{document}